\chapter*{Pronunciation \& Grammar Notes}
\addcontentsline{toc}{chapter}{Pronunciation \& Grammar Notes}
\markboth{Pronunciation \& Grammar}{Pronunciation \& Grammar}

\section*{Pronunciation Guide}

Russian pronunciation follows clear stress and palatalization rules. The transcriptions in this book use standard phonetic markers:

\begin{itemize}[leftmargin=1.5em, itemsep=0.4em]
  \item \textbf{Stress Mark} [\textquotesingle]: Indicates the stressed syllable (e.g., \textit{Мировой} {\ipafont [mira'voj]}). If no mark is present, the word is single-syllabled.
  \item \textbf{Soft Consonant} [\textquotesingle{} / ’]: Indicates palatalized (soft) consonants (e.g., \textit{Ударить} {\ipafont [u'darit’]}).
  \item \textbf{[{\ipafont ɛ}]}: Open front unrounded vowel, similar to \textit{e} in \textit{bed} or \textit{a} in \textit{cat}.
  \item \textbf{[{\ipafont ə}]}: Neutral vowel (schwa), similar to the second \textit{e} in \textit{letter}.
\end{itemize}

\section*{Grammar Notes}

Russian is an inflected language with six grammatical cases, verb aspects, and gender agreements:
\begin{itemize}[leftmargin=1.5em, itemsep=0.4em]
  \item \textbf{Aspect}: Russian verbs exist in pairs: \textit{imperfective} (ongoing, repeated actions) and \textit{perfective} (completed actions with a specific result).
  \item \textbf{Impersonal}: Constructions without an explicit grammatical subject (e.g., \textit{надо}, \textit{можно}).
  \item \textbf{Reflexive}: Verbs ending in \textit{-ся / -сь}, indicating actions performed on oneself or mutual actions.
\end{itemize}
